\chapter{Fundamentação Teórica}

Neste capítulo, são apresentados os conceitos que fundamentam a implementação de um sistema inteligente na agricultura, abrangendo os métodos e o manejo da irrigação, a Internet das Coisas aplicada à agricultura, os sensores e atuadores, os microcontroladores, os protocolos de comunicação e o armazenamento de dados.

\section{Métodos e manejo da irrigação}

\subsection{Métodos de irrigação}

Os métodos de irrigação diferem quanto à forma de aplicar a água ao solo e, consequentemente, quanto à eficiência, ao custo e à demanda de energia. Em escala global, estima-se que a substituição dos sistemas por superfície pela aspersão ou pelo gotejamento reduziria o consumo não benéfico de água, isto é, aquele perdido por evaporação, interceptação e condução, em 54\% e 76\%, respectivamente, sem redução na produtividade das culturas \cite{jagermeyr2015}.

A inundação, modalidade de irrigação por superfície usada principalmente no cultivo de arroz, ainda ocupava 44\% da área irrigada mapeada no Brasil em 2021 \cite{mapbiomas2022irrigacao}. É um método simples e com baixa demanda de energia e de materiais \cite{vigo2026}, mas exige grandes volumes de água: o arroz inundado consome de três a cinco vezes mais água que outros cereais, além de favorecer a emissão de metano \cite{ishfaq2020}.

Na aspersão, a água é aplicada sob pressão na forma de chuva artificial. O pivô central, principal sistema desse grupo, chegou a 50\% da área irrigada no Brasil em 2021 \cite{mapbiomas2022irrigacao}. O método economiza cerca de 39\% de água em relação à irrigação por gravidade \cite{chauhdary2024}, porém é sensível ao vento: as perdas por deriva e evaporação podem chegar a 25\%, com média de 11,6\% \cite{manke2019}.

O gotejamento aplica a água em baixas vazões diretamente na zona radicular. Em meta-análise, aumentou a produtividade em 28,92\% em relação à inundação, além de reduzir a lixiviação de fertilizantes e a salinidade do solo \cite{yang2023}. Suas principais desvantagens são o entupimento dos emissores, que compromete a uniformidade da irrigação \cite{shi2022}, e a maior demanda de energia e de materiais, por ser um sistema pressurizado \cite{vigo2026}.

\subsection{Manejo da irrigação}

O manejo da irrigação consiste em definir quando irrigar e quanto aplicar de água no solo, de modo a atender à necessidade da cultura sem desperdício de recursos. As principais abordagens baseiam-se no monitoramento da umidade do solo, na estimativa da evapotranspiração a partir de dados climáticos ou em indicadores da própria planta, como a temperatura do dossel \cite{datta2023}.

\subsubsection{Manejo via solo}

No manejo via solo, a decisão de irrigar se baseia na quantidade de água armazenada na zona radicular da cultura. A água disponível para as plantas corresponde à faixa entre a capacidade de campo, limite superior de retenção após a drenagem do excesso, e o ponto de murcha permanente, abaixo do qual as raízes não conseguem mais extrair água \cite{allen1998}.

Na prática, sensores instalados no solo permitem manter a umidade dentro de uma faixa ideal, cujo limite superior corresponde à capacidade de campo e o inferior ao limiar de estresse hídrico. Em um pomar de peras, essa estratégia, aplicada com uma rede de sensores sem fio, reduziu em até 50\% o volume de água em relação ao manejo convencional, mantendo a qualidade dos frutos \cite{hamouda2024}.

Em revisão sistemática, o manejo baseado em sensores de umidade proporcionou economia média de água entre 1\% e 38\% em relação a outros métodos, com produtividade igual ou superior \cite{datta2023}. A principal limitação é a precisão das leituras, que depende fortemente da calibração dos sensores para o solo em que são instalados \cite{datta2023}.

\subsubsection{Manejo via clima}

No manejo via clima, a necessidade hídrica é estimada pela \sigla{ET$_c$}{Evapotranspiração da cultura}, que corresponde à água perdida pela evaporação do solo e pela transpiração das plantas. Pelo método do boletim FAO-56, a ET$_c$ é obtida multiplicando-se a \sigla{ET$_0$}{Evapotranspiração de referência} por um \sigla{K$_c$}{Coeficiente de cultura}, que varia conforme a espécie e a fase de desenvolvimento. A ET$_0$ é calculada pela equação de Penman-Monteith, a partir de dados de temperatura e umidade do ar, velocidade do vento e radiação solar \cite{allen1998}.

A ampla adoção desse método deve-se à sua simplicidade, à robustez dos procedimentos e à possibilidade de transferir os coeficientes de cultura entre diferentes regiões \cite{pereira2015}. Por outro lado, o cálculo exige diversas variáveis meteorológicas, e o alto custo de aquisição e manutenção de estações agrometeorológicas dificulta sua aplicação, especialmente em áreas rurais \cite{moratiel2020}.

Em comparação direta, um sistema de baixo custo baseado em sensores capacitivos de umidade do solo e IoT utilizou 28,8\% menos água e alcançou produtividade da água 52,5\% maior que o manejo via clima no cultivo de alface por gotejamento. Esse resultado reforça o potencial do monitoramento do solo em tempo real, aliado à IoT, como alternativa acessível para o manejo automatizado da irrigação \cite{abdelmoneim2025}.

\section{Internet das Coisas aplicada à agricultura}

A Internet das Coisas (IoT) é um paradigma em que máquinas e objetos físicos, equipados com sensores e conectados à internet, passam a se comunicar, processar dados e coordenar ações entre si \cite{burhan2018}. Na agricultura, a tecnologia tem tornado o setor mais orientado por dados, preciso e inteligente, com o uso de sensores sem fio em etapas que vão do preparo do solo à irrigação e à colheita \cite{ayaz2019}.

Ao possibilitar o monitoramento contínuo de variáveis como umidade, temperatura e composição do solo, a IoT contribui para elevar a produção e reduzir perdas no campo \cite{rajak2023}. Entre suas principais aplicações estão os sistemas de monitoramento da irrigação, com vantagens como maior eficiência e menor uso de recursos \cite{morchid2024}.

Apesar desses benefícios, a adoção da IoT no campo ainda enfrenta limitações, como o custo de implantação, a segurança dos dados e a falta de conhecimento digital por parte dos produtores \cite{rajak2023}. Também há desafios técnicos, como o consumo de energia dos dispositivos, a latência da rede e a interoperabilidade entre as camadas do sistema \cite{villahenriksen2020}.

\subsection{Arquitetura de sistemas IoT}

Para organizar seus componentes, os sistemas IoT costumam ser descritos por arquiteturas em camadas, cujo número varia conforme o autor \cite{burhan2018}. O modelo mais difundido possui três camadas: percepção, transporte e aplicação \cite{tewari2020}.

Para a agricultura inteligente, \citeonline{morchid2024} propõem uma arquitetura de quatro camadas. A camada de percepção reúne os sensores e atuadores que coletam dados do ambiente e agem sobre ele, como sensores de umidade do solo e válvulas de irrigação. A camada de rede transmite esses dados; a camada de nuvem os armazena e processa; e a camada de aplicação apresenta as informações ao usuário, por meio de painéis ou aplicativos.

Essa organização orienta as seções seguintes deste trabalho: os sensores e microcontroladores correspondem à camada de percepção, os protocolos de comunicação à camada de rede, e o \textit{backend} com banco de dados às camadas de nuvem e de aplicação.


\section{Sensores e atuadores}

\subsection{Sensor de umidade do solo}

Um sensor de umidade do solo bastante utilizado em projetos IoT é do tipo resistivo, composto por duas hastes metálicas inseridas no solo e por um módulo eletrônico. Seu funcionamento baseia-se na condutividade elétrica do solo: quanto maior o teor de água, menor a resistência entre as hastes. O módulo converte essa variação em um sinal de tensão analógico, que é lido pelo Microcontrolador \cite{adla2020}.

Sensores resistivos têm custo muito baixo, mas são menos exatos e precisos que os capacitivos. Em laboratório, o modelo resistivo YL-69 apresentou erro (RMSE) de 5,54\% na estimativa do teor de água, contra 2,36\% a 2,88\% dos sensores capacitivos avaliados, desempenho considerado adequado para sua faixa de preço \cite{adla2020}.

A exatidão desses sensores depende da calibração para o solo e da resolução do conversor analógico digital. Usando apenas o conversor interno do Arduino, o erro de um sensor resistivo chegou a 4,79\%, caindo para 2,64\% com um conversor externo de 16 bits \cite{nandi2024}.

\subsection{Sensor de temperatura e umidade do ar}

O DHT11 é um sensor digital que mede simultaneamente a temperatura e a umidade relativa do ar. Ele combina um elemento resistivo de umidade e um termistor do tipo NTC, conectados a um microcontrolador interno de 8 bits, que envia os dados já calibrados por um barramento serial de um único fio \cite{aosongdht11}.

Segundo o fabricante, o sensor opera entre 0 e 50~$^{\circ}$C e entre 20\% e 90\% de umidade relativa, com exatidão de $\pm$2~$^{\circ}$C e $\pm$5\%, respectivamente \cite{aosongdht11}. No contexto da irrigação, essas variáveis são relevantes porque influenciam diretamente a evapotranspiração e, portanto, a demanda hídrica da cultura \cite{allen1998}.

\subsection{Atuadores: relé e válvula solenoide}

Os atuadores são os componentes que agem sobre o ambiente a partir dos comandos do microcontrolador. Na irrigação é comum encontrar válvula solenoide de 12~V, que abre ou fecha a passagem de água conforme sua bobina é energizada. Como o microcontrolador opera com baixa tensão e corrente, a válvula é comandada por meio de um relé, que funciona como uma chave eletromecânica capaz de controlar um circuito de maior potência a partir de um sinal de baixa potência \cite{ngoma2025}.

Essa configuração é comum em sistemas de irrigação automatizados. Em um sistema de gotejamento baseado no ESP32, por exemplo, o microcontrolador verifica periodicamente a umidade do solo e abre a válvula solenoide quando o solo está seco \cite{pereira2023}. A escolha do relé em sistemas de irrigação também se deve à sua facilidade de uso e à ampla disponibilidade no mercado \cite{ngoma2025}.

\subsection{Display}

Para exibir localmente as leituras dos sensores e o estado do sistema, são utilizados displays de cristal líquido (LCD) com módulo de interface I2C. Esse módulo utiliza o expansor de entradas e saídas PCF8574, que converte os comandos recebidos pelo barramento I2C em sinais paralelos para o display, reduzindo o número de pinos necessários no microcontrolador \cite{tipcf8574}.

O I2C é um barramento serial que utiliza apenas duas linhas, uma de dados (SDA) e outra de relógio (SCL), e permite conectar vários dispositivos ao mesmo barramento, com taxas de até 100~kbit/s no modo padrão e 400~kbit/s no modo rápido \cite{nxpi2c2021}.

\section{Microcontroladores e plataformas embarcadas}

Os microcontroladores e as plataformas embarcadas são responsáveis por ler os sensores, processar os dados e acionar os atuadores, compondo a camada de percepção da arquitetura IoT. Em uma revisão de 115 trabalhos sobre IoT agrícola, as plataformas abertas, como o Arduino, foram as mais utilizadas,
enquanto as soluções comerciais foram as menos adotadas, provavelmente por oferecerem pouca flexibilidade para personalização \cite{chamara2022}.

\subsection{ESP32}

O ESP32, desenvolvido pela Espressif Systems, é um \sigla{SoC}{sistema em chip} que integra um processador Xtensa LX6 de 32 bits, com até dois núcleos operando em até 240~MHz, e 520~KB de memória SRAM \cite{espressif}. Seu principal diferencial é a conectividade sem fio integrada, com Wi-Fi 802.11 b/g/n e Bluetooth, o que dispensa módulos externos de comunicação. O chip também possui conversores analógico-digitais de 12 bits com até 18 canais, permitindo a leitura direta de sensores analógicos, como o de umidade do solo \cite{espressif}.

Por esse conjunto de características, o ESP32 tem sido amplamente adotado em estações agrometeorológicas de baixo custo baseadas em IoT, sendo preferido pelo custo acessível e pela modularidade \cite{alkalaany2025}.

\subsection{Arduino}

O Arduino é uma plataforma de prototipagem baseada em hardware e software abertos. Seu modelo mais difundido, o Arduino Uno R3, utiliza o microcontrolador ATmega328P, de 8 bits, operando a 16~MHz, com 32~KB de memória flash, 2~KB de SRAM, 14 pinos digitais e 6 entradas analógicas \cite{arduinouno}.

Em sistemas de IoT agrícola descritos na literatura, o Arduino foi a placa de controle mais utilizada \cite{chamara2022}. Por outro lado, o modelo Uno não possui conectividade sem fio nativa, exigindo módulos adicionais para enviar dados à internet, e tem memória e capacidade de processamento limitadas
\cite{arduinouno}.

\subsection{Raspberry Pi}

O Raspberry Pi é um computador de placa única (\sigla{SBC}{Single-Board Computer}), capaz de executar sistemas operacionais completos, como o Linux \cite{dobrojevic2022}. O modelo Raspberry Pi 4B utiliza o processador Broadcom BCM2711, com quatro núcleos Cortex-A72 de 64 bits a 1,5~GHz, até 8~GB de memória RAM, Wi-Fi de banda dupla, Bluetooth 5.0 e Ethernet gigabit, além de um conector de 40 pinos de entrada e saída \cite{raspberrypi4}.

Esse poder de processamento permite aplicações mais complexas, como aprendizado de máquina e visão computacional. Em sistemas agrícolas, o Raspberry Pi costuma assumir o processamento dos dados, enquanto microcontroladores atuam como nós que leem os sensores e enviam as informações \cite{dobrojevic2022}. Além disso, ele não possui conversor analógico-digital integrado, o que exige um conversor externo para a leitura de sensores analógicos \cite{raspberrypi4}.

\subsection{Comparativo entre plataformas}

As três plataformas atendem a necessidades distintas. O Arduino Uno é adequado para controle local simples; o ESP32 reúne leitura de sensores e conectividade sem fio em uma única placa; e o Raspberry Pi oferece maior capacidade de processamento, mas depende de periféricos externos para leituras analógicas. A Tabela~\ref{tab:comparativo-plataformas} resume as principais características.

\begin{table}[htbp]
  \centering
  \caption{Comparativo entre as plataformas embarcadas}
  \label{tab:comparativo-plataformas}
  \small
  \begin{tabular}{p{3.2cm} p{3.4cm} p{3.4cm} p{3.4cm}}
    \hline
    \textbf{Característica} & \textbf{Arduino Uno R3} & \textbf{ESP32} & \textbf{Raspberry Pi 4B} \\
    \hline
    Tipo & Microcontrolador & SoC microcontrolado & Computador de placa única \\
    Processador & ATmega328P (8 bits) & Xtensa LX6 (32 bits, até 2 núcleos) & Cortex-A72 (64 bits, 4 núcleos) \\
    Frequência & 16 MHz & Até 240 MHz & 1,5 GHz \\
    Memória RAM & 2 KB & 520 KB & 1 a 8 GB \\
    Entradas analógicas & 6 (ADC de 10 bits) & Até 18 (ADC de 12 bits) & Não possui \\
    Conectividade sem fio & Não possui & Wi-Fi b/g/n e Bluetooth & Wi-Fi ac e Bluetooth 5.0 \\
    Sistema operacional & Não (\textit{firmware}) & Não (\textit{firmware}) & Sim (Linux) \\
    Faixa de preço\textsuperscript{*} & R\$ 28,00--57,72 & R\$ 27,33--61,90 & R\$ 495,70--920,00 \\
    \hline
  \end{tabular}
  \fonte{Elaborado pelo autor com base em \citeonline{arduinouno},
  \citeonline{espressif} e \citeonline{raspberrypi4}. \textsuperscript{*}Valores
  consultados em lojas nacionais em outubro de 2026.}
\end{table}

\section{Protocolos de comunicação}

Os protocolos de comunicação definem as regras para a troca de dados entre os dispositivos e o servidor, correspondendo à camada de rede da arquitetura IoT.

\subsection{Wi-Fi}

Wi-Fi é o nome pelo qual são conhecidas as redes locais sem fio baseadas no padrão IEEE 802.11, que especifica a camada física e a camada de controle de acesso ao meio para a conexão de dispositivos fixos, portáteis e móveis em uma área local \cite{ieee80211}. Nessas redes, os dispositivos se conectam a um ponto de acesso, geralmente um roteador, que os interliga à internet. Ao lado do LoRa, o Wi-Fi é uma das tecnologias de comunicação mais utilizadas em estações agrometeorológicas de baixo custo baseadas em IoT \cite{alkalaany2025}.

\subsection{HTTP}

O \sigla{HTTP}{Hypertext Transfer Protocol} é um protocolo da camada de aplicação, sem estado, que serve de base para a comunicação na web \cite{rfc9110}. Ele funciona no modelo requisição-resposta, o cliente envia uma requisição a um recurso identificado por uma URI, indicando um método, como GET para obter dados ou POST para enviá-los, e o servidor devolve uma resposta com um código de status, como 200 (sucesso) ou 404 (recurso não encontrado). Sua versão segura, o \sigla{HTTPS}{HyperText Transfer Protocol Secure}, transmite as mensagens de forma criptografada por meio do protocolo \sigla{TLS}{Transport Layer Security} \cite{rfc9110}.

\subsection{MQTT}

O \sigla{MQTT}{Message Queuing Telemetry Transport} é um protocolo de mensagens do tipo publicação/assinatura (\textit{publish/subscribe}), leve, aberto e simples de implementar, o que o torna adequado para ambientes com recursos limitados, como a IoT \cite{mqtt5}. Nele, os dispositivos não se comunicam diretamente: um cliente publica mensagens em um tópico, como \texttt{irrigacao/umidade}, e um servidor intermediário, chamado \textit{broker}, as repassa a todos os clientes inscritos nesse tópico. O protocolo oferece três níveis de qualidade de serviço (QoS): 0, com no máximo uma entrega; 1, com pelo menos uma entrega; e 2, com exatamente uma entrega \cite{mqtt5}.

\subsection{Comparação entre HTTP e MQTT}

No HTTP, cada envio exige uma requisição completa, com cabeçalhos em texto, o que gera uma sobrecarga elevada nas redes IoT. O MQTT, por manter uma conexão aberta com o \textit{broker} e usar cabeçalhos compactos, exige menos recursos de rede \cite{yokotani2016}. Em uma plataforma IoT baseada no módulo NodeMCU, o uso do MQTT com QoS 0 e 1 reduziu o consumo de energia em 6,03\% e 8,33\%, respectivamente, em relação ao HTTP \cite{jaraochoa2023}. Os dois protocolos também podem ser usados de forma complementar: o MQTT no envio contínuo das leituras dos sensores e o HTTP na comunicação entre o servidor e a aplicação do usuário. A Tabela~\ref{tab:http-mqtt} resume as principais diferenças.

\begin{table}[htbp]
  \centering
  \caption{Comparação entre os protocolos HTTP e MQTT}
  \label{tab:http-mqtt}
  \small
  \begin{tabular}{p{3.5cm} p{5cm} p{5cm}}
    \hline
    \textbf{Característica} & \textbf{HTTP} & \textbf{MQTT} \\
    \hline
    Modelo de comunicação & Requisição-resposta & Publicação/assinatura \\
    Intermediário & Não há (cliente e servidor) & \textit{Broker} \\
    Conexão & Uma requisição por envio & Conexão persistente \\
    Sobrecarga & Maior (cabeçalhos em texto) & Menor (cabeçalhos compactos) \\
    Qualidade de serviço & Não define níveis & Três níveis (QoS 0, 1 e 2) \\
    Uso típico & APIs web e aplicações & Telemetria de sensores \\
    \hline
  \end{tabular}
  \fonte{Elaborado pelo autor com base em \citeonline{rfc9110},
  \citeonline{mqtt5} e \citeonline{yokotani2016}.}
\end{table}

\section{Backend e armazenamento de dados}

O \textit{backend} é a parte do sistema executada no servidor, responsável por receber os dados dos dispositivos, aplicar as regras do sistema, armazenar as informações e disponibilizá-las às aplicações do usuário. Em sistemas de IoT agrícola, ele completa o fluxo que vai da coleta e transmissão dos dados ao armazenamento, à visualização, à análise e à tomada de decisão \cite{chamara2022}. 

\subsection{API}

Uma \sigla{API}{Application Programming Interface} é um conjunto de regras e pontos de acesso que permite a comunicação entre sistemas diferentes, sem que um precise conhecer a implementação interna do outro. Em aplicações web, é comum o uso do estilo arquitetural \sigla{REST}{Representational State Transfer}, proposto por \citeonline{fielding2000}, no qual cada recurso é identificado por uma URI e manipulado por meio dos métodos HTTP, em uma comunicação cliente-servidor sem estado. Os dados trocados costumam seguir o formato JSON, um padrão textual leve e independente de linguagem de programação \cite{rfc8259}.

\subsection{Banco de dados}

Um banco de dados é uma coleção de dados relacionados, organizada para atender a um propósito específico. O software que permite criá-lo, mantê-lo e consultá-lo é chamado de Sistema Gerenciador de Banco de Dados (SGBD) \cite{elmasri2018}. Nos bancos relacionais, os dados são organizados em tabelas com estrutura definida e consultados pela linguagem SQL, já os bancos NoSQL adotam modelos mais flexíveis, como documentos ou pares chave-valor, e são voltados a grandes volumes de dados e alta escalabilidade \cite{elmasri2018}. Em sistemas IoT, o banco de dados armazena o histórico das leituras dos sensores e dos acionamentos, permitindo acompanhar e analisar os dados ao longo do tempo \cite{chamara2022}. 
